\setrunninghead{序言}
\section*{序言}
\addcontentsline{toc}{section}{\protect\textbf{序言}}
A股有很多不错的股票基金品种，长期能给投资者提供比较高的回报。不过A股也存在一个奇怪的现象：大部分投资者投资股票基金的时候，很难赚到钱。很多老股民说A股“七亏二平一赚”，还是很有道理的。

一个很重要的原因，就是绝大部分投资者是在A股4000点以上，看到身边的亲戚朋友都赚钱了，才开始买股票基金的。

我们身边的亲戚朋友什么时候讨论股票最多呢？往往是在牛市最火热的阶段。他们的股票账户是什么时候开通的呢？很多人是在2007年牛市和2015年牛市时开通的。

许多投资者不知道股票品种的风险，涨得越多，买得越多。他们在牛市高位投入全部资金，之后市场下跌，收益自然不理想。这也是过去十几年，A股股票基金投资者不怎么赚钱的主要原因。

但有一种方式，可以改变投资者投资股票基金不赚钱的情况。

那就是定投。

就像“股神”沃伦·巴菲特（Warren Buffett）说的，“通过定投指数基金，一个什么都不懂的业余投资者，竟然能够战胜大部分专业投资者。”

定投是分散投资，定投的投资者不可能在牛市高位投入全部资金，风险自然大大降低。能长期坚持定投，最后大概率是赚钱的。

但是定投说起来容易，做起来难。

大部分投资者是难以在动荡的市场环境中做到坚持定投的。如何坚持下来，是投资者定投时面临的最大难点。

但这一问题，也在发生改变。比如说在基金定投方面，银行螺丝钉做得就很好。他不仅自己做好了定投，也帮助了很多人坚持长期定投，解决了投资者定投时的难点。

最早认识银行螺丝钉，是在我创办的雪球网站上。早在2014年的时候，我就看到银行螺丝钉在雪球上发表一些关于基金投资分析的文章，并坚持在收盘后发布指数基金的估值数据，回答一些网友的问题。

简单来说，他做的是基金定投，并把定投的知识和需要的数据，在网上进行分享。但就是这么一件简单的事情，他一干就是五年。

这五年里，他坚持写了近千篇文章，合计400多万字，总阅读量超过5亿人次，帮助超过百万投资者开始定投之旅。并且，他在每个交易日晚上发布估值表，无一日间断。

银行螺丝钉做的是指数基金定投，帮助大家在熊市低迷的时候定投基金，这在很大程度上解决了投资者“牛市高位买入亏钱”的痛点。买得不贵，长期坚持，赚钱也就水到渠成了。

念念不忘，必有回响。

我认为银行螺丝钉坚持做这件事情是很有意义的。

如果你是一个普通投资者，想给自己的收入找一个好的去处，如果你是一个上班族，想了解定投，我推荐你读这本书，一定会大有收获的。

\hfill {\kaishu 方三文}

\hfill {\kaishu 雪球网创始人}
